\documentclass[conference]{IEEEtran}
\IEEEoverridecommandlockouts
\usepackage{cite}
\usepackage{amsmath,amssymb,amsfonts}
\usepackage{graphicx}
\usepackage{textcomp}
\usepackage{xcolor}
\usepackage{booktabs}
\usepackage{tabularx}
\usepackage{array}
\usepackage{url}
\usepackage{balance}
\def\BibTeX{{\rm B\kern-.05em{\sc i\kern-.025em b}\kern-.08em\TeX}}

\begin{document}

\title{Why Do We Fall in Love?\\
Neuroendocrine Pair-Bonding Causal Atlas}

\author{\IEEEauthorblockN{Source-Controlled Research Plan}
\IEEEauthorblockA{Survey--Idea--ExperimentDesign--Author pipeline\\
Neuroscience and biological mechanisms only}}

\maketitle

\begin{abstract}
Romantic love is often described with a short list of hormones, but a list does not explain how a transient sexual motivation can become selective partner-directed reward and then a persistent pair bond. This paper proposes a Neuroendocrine Pair-Bonding Causal Atlas (NPB-CA), a staged cross-species research program that treats these phenomena as interacting biological states rather than as a single molecule or anatomical center. The causal arm uses prairie voles to perturb oxytocin-receptor, vasopressin-receptor, nucleus-accumbens dopamine, and ventral-tegmental-area-to-nucleus-accumbens pathways during bond formation and maintenance. The translational arm uses repeated, non-invasive human neuroimaging and endocrine measurements with partner, familiar non-partner, stranger, and nonsocial controls. A phase-indexed state-space model and shared feature-space representational analysis separate partner selectivity from locomotion, general social approach, arousal, nonsocial reward, and stress. The proposal makes direct, falsifiable predictions about receptor specificity, D1/D2 phase interactions, dopamine-linked partner cues, endocrine moderation, and cross-session neural stability. It does not claim that animal behavior is identical to human experience, and it reports no newly observed results. Its central scientific claim is prospective: falling in love is most plausibly produced by a transition in a distributed neuroendocrine circuit in which receptor-localized social recognition gates mesolimbic reward and endocrine state modulates the gain and persistence of the transition.
\end{abstract}

\begin{IEEEkeywords}
pair bonding, romantic love, oxytocin, vasopressin, dopamine, nucleus accumbens, ventral tegmental area, neuroendocrinology, prairie vole, causal atlas
\end{IEEEkeywords}

\section{Introduction}
The phrase ``falling in love'' compresses several biological events into one label. Sexual motivation, selection of one partner from many possible social targets, reinforcement of partner-associated cues, and long-term maintenance are related but experimentally separable. The useful question is therefore not where a single emotion is located. It is how endocrine state, neuromodulator release, receptor distribution, synaptic plasticity, and distributed neural circuits interact across time to produce a selective and persistent pattern of behavior.

The distinction matters for both mechanism and translation. A manipulation can reduce social contact because it impairs locomotion, arousal, stress regulation, or sensory processing without altering partner selectivity. Conversely, an intervention can alter preference formation but have little effect once a bond is established. A biological account must therefore measure phase, cue specificity, and non-target functions together.

This paper presents the NPB-CA research plan generated through a four-stage evidence-to-design workflow. The Survey stage maps evidence and gaps; the Idea stage selects an integrated causal direction; ExperimentDesign converts the direction into a preregisterable protocol; and the Author stage binds every proposed claim to the upstream source registry. The scope is intentionally limited to biology and neuroscience. The report does not analyze love through philosophy, culture, literary interpretation, or general social-science constructs. It also does not report a completed animal or human experiment.

The central proposal is that partner bonding is a state transition in a distributed circuit. Gonadal-axis and stress-axis variables change the gain of hypothalamic and mesolimbic processing. Oxytocin and vasopressin receptor fields regulate the selectivity of partner cues in context-dependent regions. VTA--NAc dopamine assigns incentive value during formation and contributes differently during maintenance. Repeated interaction and plasticity stabilize the resulting partner preference. The strongest version of this model is falsifiable: it predicts phase-specific and partner-specific effects of perturbation that cannot be explained by general social contact or arousal.

\section{Background, Survey, and Research Gap}
\subsection{Pair bonding as a measurable biological program}
Prairie voles provide a valuable causal model because a defined mating or cohabitation history can be followed by partner preference and bond-maintenance assays. Young and Wang synthesized evidence that pair bonding involves coordinated oxytocin, vasopressin, and dopamine mechanisms rather than an isolated signal \cite{young2004}. Insel and Young similarly emphasized receptor localization and neural circuit context in the neurobiology of attachment \cite{insel2001}. These studies justify a mechanistic program, but they do not warrant the claim that vole partner preference is identical to human romantic love. The valid bridge is through measurable operations: cue discrimination, reward-linked neural activity, selective approach, and persistence after delay.

\subsection{Dopamine and the formation--maintenance distinction}
The mesolimbic system is a candidate substrate for the conversion of a partner cue into a reinforced biological target. In prairie voles, NAc dopamine has different contributions to formation and maintenance of monogamous pair bonds, with D1- and D2-family receptors playing non-equivalent roles \cite{aragona2006}. This result is more informative than the popular claim that dopamine is simply a happiness chemical. Dopamine is treated here as a circuit signal that changes cue value, prediction, and motivated approach. The design directly tests whether receptor effects are phase-dependent and whether they remain after controlling for general movement and social investigation.

\subsection{Neuropeptide receptor fields}
Oxytocin and arginine vasopressin act through receptors whose expression is distributed and regionally heterogeneous. A receptor's effect depends on where it is expressed, which cells carry it, the current endocrine state, and the sensory context. Later comparative work continues to use voles to study neuropeptidergic regulation, stress buffering, and pair-bonding mechanisms \cite{johnson2015}. Oxytocin-related findings also include social consolation and developmental effects, showing why general social contact and stress physiology must be included as negative controls rather than assumed to be equivalent to bonding \cite{bales2023,burkett2016}.

\subsection{Human evidence and its limits}
Human neuroimaging has reported activity in reward and motivation-related regions during exposure to a beloved partner \cite{bartels2000}. A comparative review framed romantic love as a mammalian brain system for mate choice and connected human imaging to animal neurobiology \cite{fisher2006}. These data provide translational motivation, not causal proof. Partner images differ from unfamiliar images in familiarity, visual salience, autobiographical association, arousal, and learned reward. NPB-CA therefore uses partner minus familiar non-partner as its primary human contrast, with stranger and nonsocial conditions to test specificity. The predicted output is a distributed, reproducible neural representation, not a ``love center.''

\subsection{Endocrine modulation and genetic variation}
Testosterone and estradiol are primarily produced by gonadal and adrenal tissues; the hypothalamus and pituitary coordinate their regulation. They are modeled as state variables that alter circuit gain, not as substances produced exclusively in the hypothalamus. Corticosterone or cortisol is included because stress-axis activity can change social investigation, reward learning, and peptide release. The AVPR1A association reported by Walum et al. connects receptor-related genetic variation to pair-bonding measures in humans \cite{walum2008}. It is a candidate moderator, not a deterministic ``love gene.'' Replication, ancestry-aware modeling, sex, age, medication, and endocrine covariates are required before any general inference.

\subsection{Gap ledger}
The survey produced six accepted gaps. G1 is the lack of unified phase boundaries for sexual motivation, partner-directed reward, and durable preference. G2 is the absence of a common causal decomposition of VTA--NAc dopamine, hypothalamic neuropeptides, and limbic recognition circuits. G3 concerns the temporal and interaction logic of dopamine, oxytocin, vasopressin, sex steroids, and stress hormones. G4 is the missing measurement language between human imaging and vole perturbation. G5 is the specificity problem: partner selectivity must be separated from general affiliation, locomotion, novelty, sexual behavior, and stress buffering. G6 is the need for adequately replicated genetic and endocrine moderation analyses.

\begin{table*}[t]
\caption{Evidence anchors and permitted inference}
\label{tab:evidence}
\centering
\small
\begin{tabularx}{\textwidth}{p{0.08\textwidth}p{0.29\textwidth}p{0.28\textwidth}X}
\toprule
ID & Anchor & Biological contribution & Inference boundary \\
\midrule
E1 & Young and Wang, 2004 \cite{young2004} & Pair-bonding neurobiology and vole circuit synthesis & Direct mechanistic background; species-bounded \\
E2 & Insel and Young, 2001 \cite{insel2001} & Attachment-related receptor and circuit organization & Supports distributed architecture, not a single center \\
E3 & Bartels and Zeki, 2000 \cite{bartels2000} & Human partner-cue neuroimaging & Correlational human evidence \\
E4 & Fisher, Aron, and Brown, 2006 \cite{fisher2006} & Comparative mate-choice brain-system synthesis & Comparative hypothesis, not causal human evidence \\
E5 & Aragona et al., 2006 \cite{aragona2006} & NAc dopamine during formation and maintenance & Strong animal causal anchor \\
E6 & Johnson and Young, 2015 \cite{johnson2015} & Social attachment and pair-bond circuit mechanisms & Mechanistic synthesis with species limits \\
E7 & Walum et al., 2008 \cite{walum2008} & AVPR1A variation and human pair-bonding measures & Association; requires replication \\
E8 & Bales, 2023; Burkett et al., 2016 \cite{bales2023,burkett2016} & Oxytocin development, social behavior, and consolation & Supports context and specificity controls \\
\bottomrule
\end{tabularx}
\end{table*}

\section{Research Questions and Planned Contributions}
NPB-CA asks three linked questions. First, which receptor and mesolimbic circuit perturbations selectively alter partner preference during formation versus maintenance? Second, how do dopamine, oxytocin, vasopressin, gonadal-axis variables, and stress-axis variables interact over the bond trajectory? Third, can a prespecified neural representation be compared across prairie voles and humans without equating their behavioral states?

The planned contributions are fourfold. (1) A phase-indexed operational vocabulary separates formation from maintenance. (2) A causal perturbation matrix tests receptor and dopamine contributions with non-target controls. (3) A shared feature dictionary makes cross-species comparison possible without assuming anatomical identity. (4) A failure-aware analysis plan defines what would reject a strong conserved-mechanism claim.

The proposal is deliberately more ambitious than a descriptive hormone survey. It seeks a causal atlas that predicts when partner cues acquire selectivity, how that selectivity is stabilized, and which biological variables explain its gain. A positive result would organize previously separate findings into a testable dynamic model. A negative result would still identify which popular simplifications fail: a single hormone, a single brain site, or a universal receptor effect.

\section{Idea Source Checkpoints and Direction Selection Audit}
The Idea stage began with three routes. A single-hormone explanation was rejected because it cannot encode receptor localization, phase transitions, or nonspecific effects. A human-imaging-only route was retained as a translational component but rejected as a complete causal program. A vole-only receptor atlas was retained as the causal anchor but could not address human neural representation. NPB-CA combines the latter two with endocrine time series and a common feature space.

The direction was selected by five criteria: alignment to all six accepted gaps, explicit falsifiability, causal leverage, cross-species value, and staged feasibility. Its highest-value design change was replacing a static ``love signature'' with formation and maintenance states. Its second was replacing anatomical equivalence with representational geometry. Its third was adding nonsocial reward, locomotion, and stress controls so that a partner-preference effect has a specific biological interpretation.

\begin{table}[t]
\caption{Direction selection audit}
\label{tab:audit}
\centering
\small
\begin{tabular}{lccccc}
\toprule
Route & Novelty & Causal & Translation & Gaps & Decision \\
\midrule
Single hormone & Low & Low & Medium & 2/6 & Reject \\
Human imaging only & Medium & Low & High & 3/6 & Component \\
Vole receptor only & High & High & Medium & 4/6 & Component \\
NPB-CA & High & High & High & 6/6 & Primary \\
\bottomrule
\end{tabular}
\end{table}

\section{Problem Definition, Assumptions, and Hypotheses}
For this study, a bond is a persistent, partner-selective biological phenotype measured by a partner-versus-stranger preference index after standardized exposure and delay. It is not defined by a self-report scale and is not claimed to exhaust the human phenomenon. Sexual motivation is indexed by reproductive-state and mating-related measures; partner-directed reward is indexed by cue-selective dopamine and neural responses; maintenance is indexed by delayed preference persistence and selective contact.

Let the latent state for individual (i) at time (t) be
\begin{equation}
\mathbf{x}_{it}=[d_{it},r_{it},q_{it},s_{it},g_{it},p_{it}]^{\mathsf T},
\end{equation}
where (d) is dopamine-linked cue value, (r) receptor-gated selectivity, (q) persistence/plasticity, (s) stress-axis state, (g) gonadal-axis state, and (p) peptide-related signaling. In phase (phi),
\begin{equation}
\mathbf{x}_{i,t+1}=\mathbf{A}_{\phi}\mathbf{x}_{it}+\mathbf{B}_{\phi}\mathbf{u}_{it}+\mathbf{G}_{\phi}\mathbf{h}_{it}+\boldsymbol{\epsilon}_{it}.
\end{equation}
The observation model is
\begin{equation}
\mathbf{y}_{it}=f_{\phi}(\mathbf{C}_{\phi}\mathbf{x}_{it}+\mathbf{D}\mathbf{c}_{it})+\boldsymbol{\eta}_{it},
\end{equation}
where $\mathbf{h}$ contains measured endocrine variables and $\mathbf{c}$ contains cue class, exposure, locomotion, motion, and assay-batch covariates.

The causal estimand is
\begin{equation}
\Delta_{k,m,\phi}=E[Y_m\mid do(U_k=1),\phi]-E[Y_m\mid do(U_k=0),\phi],
\end{equation}
where $k$ is a perturbation target and $m$ is partner selectivity, general social approach, nonsocial reward, locomotion, stress physiology, or persistence.

The hypotheses are:
\textbf{H1}: OXTR or AVPR1A perturbation selectively reduces partner preference after locomotion and general social-contact controls.\\
\textbf{H2}: NAc D1 and D2 perturbations produce a formation-by-maintenance interaction.\\
\textbf{H3}: VTA--NAc dopamine responses are more partner-selective than matched stranger and nonsocial-reward responses during formation.\\
\textbf{H4}: endocrine and stress-axis state moderates but does not replace circuit-level prediction.\\
\textbf{H5}: human partner-cue neural representations replicate across sessions and remain distinct from familiar-person and nonsocial representations.\\
\textbf{H6}: a low-dimensional feature geometry is directionally shared across species, with nonzero species-specific deviations.

\section{Study Design and Methods}
\subsection{Causal animal arm}
The planning sample is 96 adult prairie voles, balanced by sex when colony availability permits, randomized within litter and housing block. This is a planning target subject to pilot-based power calculation and animal-care approval, not an observed sample. Phases are baseline, formation, established maintenance, and delayed retention. Cohabitation, handling, arena exposure, light cycle, and food access are standardized. A non-bonded repeated-social-exposure condition separates partner-specific learning from social exposure alone.

The perturbation matrix includes regionally restricted OXTR signaling, AVPR1A signaling, NAc D1-family signaling, NAc D2-family signaling, and VTA--NAc dopamine activity, plus vehicle, sham, or pathway-control groups. The exact validated pharmacological, chemogenetic, or optical implementation is chosen after local technical and welfare review. Positive findings require an orthogonal perturbation where feasible, reducing the chance that a drug or actuator artifact is interpreted as a mechanism.

The primary behavioral endpoint is the partner-selectivity index
\begin{equation}
PSI=\log\left(\frac{T_{partner}+\epsilon}{T_{stranger}+\epsilon}\right).
\end{equation}
Secondary endpoints include total social contact, stranger contact, locomotion, object investigation, naturally occurring mating-related behavior, and delayed persistence. A partner-specific mechanism must alter PSI more than locomotion or total social contact and must not be explained by a change in object investigation or nonsocial reward.

Neural measures include NAc dopamine fluorescence or electrophysiology during partner, stranger, and nonsocial reward cues; VTA activity where technically validated; and receptor expression or occupancy in predefined regions. Endocrine samples include oxytocin, vasopressin where assay performance is acceptable, testosterone or estradiol as sex-specific covariates, and corticosterone. Sampling time, handling, and assay batch are standardized. Because peripheral peptide measurements can be noisy and may not represent local synaptic release, they are treated as covariates rather than direct proxies for circuit concentration.

\subsection{Human translational arm}
The planning sample is 120 adults in established pair relationships, with balanced sex representation where recruitment permits. Each participant completes two sessions separated by at least two weeks, with screening for acute illness, severe sleep loss, and medication-related confounds. The arm collects non-invasive fMRI, pulse, respiration, skin conductance, and saliva or blood endocrine covariates under institutional review. It includes no hormone administration, brain stimulation, deception, or relationship intervention.

Four cue classes are prespecified: current partner, familiar non-partner, unfamiliar person, and nonsocial object or reward. Low-level visual properties are matched, and order is counterbalanced. The primary contrast is partner minus familiar non-partner; partner minus stranger assesses broad selectivity; partner minus nonsocial reward tests social specificity. Repeated partner cues estimate within-participant stability. A brief participant-provided class label can document stimulus identity, but the biological endpoints are neural, physiological, and endocrine measures.

The primary imaging regions are VTA, NAc or ventral striatum, hypothalamus, amygdala, hippocampus, septal and ventral-pallidal regions where resolution permits, and medial prefrontal cortex. Whole-brain analysis is secondary and multiplicity-corrected. Cross-validated multivoxel pattern analysis tests whether partner-versus-control representations generalize across sessions. Endocrine variables moderate imaging estimates; they are not interpreted as direct readouts of love.

\subsection{Cross-species feature space}
The bridge uses a shared feature dictionary rather than anatomical identity. Each species contributes partner-cue selectivity, stranger selectivity, nonsocial-reward response, locomotion or arousal control, delayed persistence, dopamine-linked cue response, and endocrine state. Rank-based representational similarity compares the arrangement of partner, familiar, stranger, and object conditions. A hierarchical model is
\begin{equation}
\theta_{s,k}=\mu_k+\delta_{s,k},\quad \delta_{s,k}\sim N(0,\tau_k^2),
\end{equation}
where (s) indexes species and (k) indexes shared features. The bridge is accepted only if the shared effect is directionally consistent and out-of-sample prediction exceeds a preregistered baseline. A failure rejects the strong conservation claim but preserves species-specific mechanisms.

\subsection{Statistical analysis and reproducibility}
The animal primary analysis is a mixed-effects model with phase, perturbation, cue class, sex, endocrine covariates, and prespecified interactions; animal and housing block are random effects. The human analysis combines cross-validated decoding with a hierarchical region-level model. Missing data, assay limits, motion thresholds, and exclusion rules are declared before unblinding. Endpoint families are corrected for multiplicity.

Power is calculated separately for the animal causal contrast and human cross-session decoding from pilot variance, minimum biologically meaningful effect, and expected attrition. The planning sample numbers are not evidence. Code, stimulus metadata, preregistration, perturbation verification, exclusion decisions, and model specifications are versioned. The primary atlas claim requires convergence across perturbation, neural readout, phase interaction, and specificity controls; a single significant biomarker is insufficient.

\subsection{Model comparison and quality-control hierarchy}
The state-space model is not assumed to be correct merely because it is biologically plausible. Three nested models are compared. The null model contains cue, phase, and nuisance effects but no circuit-by-endocrine interaction. The additive model adds dopamine, peptide, gonadal-axis, and stress-axis predictors without interactions. The atlas model adds phase-specific circuit coefficients, receptor perturbation effects, and a shared latent feature layer. Predictive comparison is performed with blocked cross-validation: all observations from an animal, participant, or session remain in one fold. A more complex model is retained only when it improves held-out prediction and its added coefficients have interpretable uncertainty.

The quality-control hierarchy follows the causal question. First, perturbation verification confirms regional targeting, receptor modulation, or pathway modulation before behavioral interpretation. Second, sensory and motor controls test whether cue discrimination can be measured independently of vision, locomotion, or general arousal. Third, phase controls test whether the same intervention has a distinct formation and maintenance effect. Fourth, replication controls test whether the effect survives an orthogonal method, a new housing block, or a second human imaging session. This hierarchy avoids converting a technically impressive but nonspecific signal into a biological conclusion.

For the animal arm, the primary decision boundary is a perturbation-by-phase effect on PSI accompanied by a smaller or null effect on locomotion and total social contact. A Bayesian posterior probability or a preregistered confidence interval can express this boundary, but the numerical threshold is set from pilot variance before unblinding. For the human arm, the primary decision boundary is cross-session decoding above a preregistered null obtained by label permutation within participant and session. The cross-species boundary is out-of-sample directional agreement in feature geometry, not a requirement that human and vole effect magnitudes be equal.

Several alternative explanations are explicitly modeled. Familiarity may mimic partner salience; the familiar non-partner condition addresses this. Novelty may mimic stranger rejection; repeated exposure and counterbalancing address this. General reward may mimic NAc dopamine activity; nonsocial reward addresses this. Stress relief may mimic oxytocin-linked contact; corticosterone or cortisol and stress-linked physiology address this. Sexual motivation may change with reproductive state; sex-steroid measures and phase-specific covariates address this. Finally, a genotype association may reflect population structure; ancestry-aware covariates, cohort replication, and sensitivity analyses address this.

The data product is an auditable atlas rather than a single score. For every candidate mechanism, it records the target, phase, cue contrast, neural measure, endocrine moderator, specificity controls, causal estimate, uncertainty, replication status, and species scope. A mechanism is promoted from candidate to supported only when the corresponding row satisfies its prespecified decision rule. This format lets a later study replace a weak assay, add a species, or revise one latent transition without silently rewriting the entire biological interpretation.

\section{Expected Outcome Branches and Conditional Conclusions}
The proposal predicts five informative branches. In Branch A, receptor perturbation selectively reduces PSI while preserving locomotion and general social approach. This would support a causal role for receptor-gated partner recognition. In Branch B, D1/D2 perturbations differ between formation and maintenance. This would support phase-specific dopamine computation. In Branch C, human partner-cue representations replicate across sessions and share feature geometry with vole data. This would support a constrained translational bridge, not anatomical identity.

In Branch D, the effects are nonspecific: locomotion, stress, or general social contact changes more than PSI. The corresponding mechanistic claim is rejected. In Branch E, the cross-species bridge fails or human decoding does not generalize across sessions. The atlas then separates species-specific mechanisms and rejects the strong conserved-representation claim. These are design-time branches; no branch is presented as an observed result.

The most important conditional conclusion is therefore precise. If H1--H4 survive specificity controls and H5--H6 replicate out of sample, falling in love can be modeled as a distributed transition in which social recognition gates reward and endocrine state controls gain and persistence. If only the animal arm succeeds, the conclusion is a causal vole pair-bond mechanism with limited human translation. If only human imaging succeeds, the conclusion is a reproducible partner-cue neural representation without demonstrated causal mechanism. If neither succeeds, the single-molecule and single-center explanations lose support and the operational definitions require revision.

\section{Risks, Limitations, and Human Review Requirements}
The first risk is construct overreach. A partner-preference assay is a powerful behavioral operation but is not identical to human romantic experience. The paper avoids that equivalence by treating human and vole observations as different measurements of partly shared circuit computations. The second risk is peripheral endocrine interpretation. Circulating oxytocin or vasopressin can be unstable and may not reflect local release; assay validation and cautious moderation models are essential.

The third risk is perturbation specificity. Receptor antagonists can have off-target actions, viral or optical tools can alter tissue, and D1/D2 labels do not perfectly partition circuit function. Orthogonal perturbations, pathway controls, receptor verification, and blinded scoring reduce this risk. The fourth risk is imaging confounding. Partner cues carry visual and autobiographical differences. Matched familiar, stranger, object, physiological, and motion controls are therefore part of the primary design, not optional supplements.

The fifth risk is population heterogeneity. Sex, age, reproductive state, endocrine cycle, medication, genotype, stress, sleep, and relationship duration can alter measurements. These variables are prespecified moderators and not post hoc explanations. AVPR1A or OXTR variants must not be used to label individuals or predict relationship outcomes; any optional genotype analysis requires privacy-preserving governance and multi-cohort replication.

Animal work requires institutional animal-care approval, 3R justification, humane endpoints, minimized separation duration, and technical review of any viral, pharmacological, or optical method. Human work requires institutional review, informed consent, MRI safety screening, data minimization, and secure handling of endocrine and optional genetic data. No clinical treatment, relationship manipulation, or behavioral advice is proposed. The final judgment of biological validity remains with domain experts who inspect assay performance, perturbation verification, species interpretation, and statistical assumptions.

\section{Conclusion}
The evidence supports a direct and testable answer to why we fall in love: not because one molecule turns on a universal love center, but because endocrine state and receptor-localized social-recognition circuits interact with mesolimbic reward and plasticity to make one partner biologically salient and behaviorally persistent. That answer is a mechanistic hypothesis, not a completed discovery. NPB-CA turns it into a causal program by separating formation from maintenance, partner selectivity from general social behavior, and cross-species similarity from anatomical equivalence. Its value will be determined by preregistered perturbations, matched controls, replication across sessions and cohorts, and its willingness to treat failed branches as evidence against over-simple explanations.

\appendices
\section{Variables and Operational Definitions Appendix}
\begin{table*}[h]
\caption{Operational variables inherited by every stage}
\label{tab:variables}
\centering
\small
\begin{tabularx}{\textwidth}{p{0.2\textwidth}p{0.28\textwidth}p{0.42\textwidth}}
\toprule
Variable & Operational definition & Interpretation rule \\
\midrule
Formation & Pre-association baseline through standardized cohabitation and immediate preference assay & Tests acquisition of partner selectivity \\
Maintenance & Established-bond assay after bond confirmation and delay & Tests persistence and stabilization \\
PSI & Log ratio of partner to stranger contact time with fixed epsilon & Positive value means partner-selective contact; not a human-experience score \\
Partner-cue response & Neural or dopamine-linked response to current partner cue minus matched control & Must survive visual, motion, and arousal controls \\
General social approach & Total contact with social stimuli across partner and stranger conditions & Nonspecific social effect control \\
Nonsocial reward & Neural/behavioral response to matched object or reward cue & Separates social bonding from general reward \\
Endocrine state & Oxytocin, vasopressin where valid, testosterone/estradiol, and corticosterone/cortisol measurements & Moderator or covariate; not a single-cause label \\
Persistence & Partner selectivity retained after a prespecified delay & Evidence for maintenance, not merely immediate approach \\
Cross-species bridge & Out-of-sample agreement in shared feature geometry & Does not require anatomical or behavioral identity \\
\bottomrule
\end{tabularx}
\end{table*}

\section{Evidence Coverage and Review Appendix}
The Survey evidence plan assigns E1, E2, E5, and E6 to the animal circuit mechanism; E3 and E4 to human partner-cue imaging; E7 to candidate genetic moderation; and E8 to oxytocin context and specificity. Claims in this paper are limited to those anchors or are explicitly marked as proposed tests. The accepted gaps are G1 through G6 as listed in Section II. The primary direction is linked to every accepted gap. The document status is \texttt{PROPOSAL\_NO\_OBSERVED\_RESULTS}; the observed-results list is empty.

\balance
\begin{thebibliography}{00}
\bibitem{young2004} L. J. Young and Z. Wang, ``The neurobiology of pair bonding,'' \emph{Nature Neuroscience}, vol. 7, no. 10, pp. 1048--1054, 2004, PMID: 15452576.
\bibitem{insel2001} T. R. Insel and L. J. Young, ``The neurobiology of attachment,'' \emph{Nature Reviews Neuroscience}, vol. 2, no. 2, pp. 129--136, 2001, PMID: 11252992.
\bibitem{bartels2000} A. Bartels and S. Zeki, ``The neural basis of romantic love,'' \emph{NeuroReport}, vol. 11, no. 17, pp. 3829--3834, 2000, PMID: 11117499.
\bibitem{fisher2006} H. E. Fisher, A. Aron, and L. L. Brown, ``Romantic love: A mammalian brain system for mate choice,'' \emph{Philosophical Transactions of the Royal Society B}, vol. 361, no. 1476, pp. 2173--2186, 2006, PMID: 17118931.
\bibitem{aragona2006} B. J. Aragona, Y. Liu, Y. J. Yu, J. T. Curtis, J. M. Detwiler, T. R. Insel, and Z. Wang, ``Nucleus accumbens dopamine differentially mediates the formation and maintenance of monogamous pair bonds,'' \emph{Nature Neuroscience}, vol. 9, no. 1, pp. 133--139, 2006, PMID: 16327783.
\bibitem{johnson2015} Z. V. Johnson and L. J. Young, ``Neurobiological mechanisms of social attachment and pair bonding,'' \emph{Current Opinion in Behavioral Sciences}, vol. 3, pp. 38--44, 2015, PMID: 26146650.
\bibitem{walum2008} H. Walum et al., ``Genetic variation in the vasopressin receptor 1a gene (AVPR1A) associates with pair-bonding behavior in humans,'' \emph{Proceedings of the National Academy of Sciences}, vol. 105, no. 37, pp. 14153--14156, 2008, PMID: 18765804.
\bibitem{bales2023} K. L. Bales, ``Oxytocin: A developmental journey,'' \emph{Comprehensive Psychoneuroendocrinology}, vol. 16, p. 100203, 2023, PMID: 38108037.
\bibitem{burkett2016} J. P. Burkett, E. Andari, Z. V. Johnson, D. C. Curry, F. B. de Waal, and L. J. Young, ``Oxytocin-dependent consolation behavior in rodents,'' \emph{Science}, vol. 351, no. 6271, pp. 375--378, 2016, PMID: 26798013, PMCID: PMC4737486.
\end{thebibliography}

\end{document}
